% Recorrido del educador: de la instalación al laboratorio en curso.
\begin{tikzpicture}[
    font=\small,
    paso/.style={
        draw, rounded corners=2pt, align=center,
        minimum height=1.45cm, text width=3.05cm, inner sep=4pt
    },
    fase/.style={font=\small\bfseries, anchor=west},
    flecha/.style={-{Stealth[length=2.2mm]}, thick}
]

% ---- Fase 1: antes de la clase -------------------------------------------
\node[fase] at (-0.7,0) {Antes de la clase};

\node[paso] (p1) at (0.9,-1.45)  {1. Instalar\\\footnotesize Docker rootless\\\footnotesize y \texttt{uv sync}};
\node[paso, right=0.5cm of p1] (p2) {2. Imagen vulnerable\\\footnotesize construir o adaptar};
\node[paso, right=0.5cm of p2] (p3) {3. Escenario en YAML\\\footnotesize objetivos y pistas};
\node[paso, right=0.5cm of p3] (p4) {4. \texttt{scenario load}\\\footnotesize y \texttt{scenario validate}};

\draw[flecha] (p1) -- (p2);
\draw[flecha] (p2) -- (p3);
\draw[flecha] (p3) -- (p4);

% ---- Fase 2: la clase ----------------------------------------------------
\node[fase] at (-0.7,-3.05) {Durante la clase};

\node[paso] (p5) at (0.9,-4.5)  {5. \texttt{daemon start}\\\footnotesize en segundo plano};
\node[paso, right=0.5cm of p5] (p6) {6. \texttt{run batch-start}\\\footnotesize un entorno por participante};
\node[paso, right=0.5cm of p6] (p7) {7. \texttt{dashboard} y \texttt{feed}\\\footnotesize avance en vivo};
\node[paso, right=0.5cm of p7] (p8) {8. \texttt{run show}\\\footnotesize evidencia y tiempos};

\draw[flecha] (p5) -- (p6);
\draw[flecha] (p6) -- (p7);
\draw[flecha] (p7) -- (p8);

% ---- Fase 3: cierre ------------------------------------------------------
\node[fase] at (-0.7,-6.1) {Al cerrar};

\node[paso, text width=4.5cm, anchor=west] (p9) at (p1.west |- 0,-7.4) {9. \texttt{scenario/participant reset-all}\\\footnotesize libera los entornos};

\end{tikzpicture}
